\section{Usuarios o clientes potenciales y contexto de uso}

Los resultados podrían interesar a las instituciones que recopilan o utilizan información de seguridad vial. No se consultó directamente a ninguna de ellas, de modo que sus intereses se plantean como hipótesis. Lo que sí está documentado es el papel de cada una en los datos, y es ese papel el que condiciona el problema.

\begin{table}[htbp]
\centering
\small
\caption{Actores, intereses hipotéticos y papel en los datos}
\label{tab:actores}
\begin{tabularx}{\textwidth}{p{2.5cm}p{3.6cm}p{3.7cm}X}
\toprule
Actor & Interés (hipótesis) & Papel en los datos & Cómo condiciona el problema \\
\midrule
ANT & Complemento a sus estadísticas & Publica el archivo y define cómo se cuentan las víctimas & Fija la definición de fallecido y las variables disponibles; el proyecto no las modifica \\
Entes de control (Policía Nacional, CTE, entidades municipales) & Referencia sobre dónde se concentran los siniestros graves & Levantan el parte en el sitio & La calidad y la cobertura dependen de su registro; las zonas con más control pueden aparecer sobrerrepresentadas \\
Gobiernos autónomos descentralizados & Información contextual para estudios de seguridad vial & Usuarios finales de resultados agregados & Solo les sirve localmente si hay suficientes casos por territorio \\
INEC & Referencia de contraste & Procesa los mismos registros para sus boletines \cite{ref1} & Permite verificar los totales \\
Equipo del proyecto & Responder la pregunta de investigación & Construye la etiqueta y la evaluación & Limitado por el tiempo, los equipos personales y las reglas del curso \\
\bottomrule
\end{tabularx}
\end{table}

Dado este conjunto de actores, el contexto de uso es el análisis posterior al registro: resultados agregados que complementan las estadísticas existentes y un prototipo académico de consulta. Ninguno de los dos se usa para anticipar siniestros, sancionar conductores ni decidir sobre territorios. La distinción entre el registro y el análisis aparece con más precisión en los modelos de la sección~5.
